\section{Introduction}\label{sec:introduction}

Let $d\ge 1$ and let $H^{(N)}:=(H_1^{(N)},\ldots,H_d^{(N)})$ consist of independent $N\times N$ matrices from the Gaussian unitary ensemble (GUE), each with density proportional to $\exp(-N\Tr(H^2)/2)$ on the real vector space of Hermitian matrices.
The tuples for different $N$ are defined on a common probability space with arbitrary coupling across sizes.
Here $\Tr$ denotes the unnormalized matrix trace and $\tr_N:=N^{-1}\Tr$.
Let $s:=(s_1,\ldots,s_d)$ be a free centered semicircular family of variance one in a finite von Neumann algebra $\mathcal M$ with faithful normal tracial state $\tau$.
Each $s_j$ has spectral density $(2\pi)^{-1}\sqrt{4-x^2}$ on $[-2,2]$.
Freeness means that the trace of a product of centered elements from the unital algebras generated by the individual $s_j$ vanishes whenever consecutive factors come from distinct algebras \cite[Definition~2.16, Remark~2.17(2), and Definition~5.3]{NS06}.

For each scalar polynomial $Q$ in the input variables, $\tr_N(Q(H^{(N)}))$ converges in probability to $\tau(Q(s))$.
The maximum of the input operator norms converges in probability to $2$, the right endpoint of the variance-one semicircle law \cite[pp.~1--2 of arXiv v1]{CGP22}.
For a self-adjoint scalar polynomial $P$, this gives convergence of the empirical eigenvalue distribution of $P(H^{(N)})$ to the spectral distribution of $P(s)$.
For a general polynomial, the eigenvalues can lie in the complex plane, and the corresponding limit is given by the Brown measure.

We use the Brown measure convention of \cite[Chapter~11, Theorem~5, Definition~6 and pp.~268--270]{MS17}.

\begin{definition}\label{def:brown-measure}
Let $\mathcal M$ be a finite von Neumann algebra with faithful normal tracial state $\tau$, and let $a\in\mathcal M$.
The Brown measure of $a$ is the measure $\mu_a^{\mathrm{Br}}$ given by
\[
 \mu_a^{\mathrm{Br}}:=\frac{1}{2\pi}\Delta u_a,\qquad
 u_a(z):=\tau(\log|a-z1_{\mathcal M}|).
\]
Here $|b|:=(b^*b)^{1/2}$, the trace of the logarithm is the spectral integral with respect to $\tau$, and $\Delta$ acts in the sense of distributions on $\mathbb C$.
\end{definition}

The potential $u_a$ is locally integrable, and $\mu_a^{\mathrm{Br}}$ is a probability measure supported in the spectrum of $a$.
Moreover, $u_a(z)=\int_{\mathbb C}\log|z-w|\,d\mu_a^{\mathrm{Br}}(w)$, and for normal $a$ the Brown measure agrees with its spectral distribution.

Belinschi, Śniady and Speicher \cite[Section~7]{BSS18} formulate the eigenvalue convergence conjecture for general polynomials in independent self-adjoint Gaussian matrices.
Han \cite[Theorem~1.2]{Han26} proves weak convergence in probability for polynomials in independent complex Ginibre matrices using only the unstarred variables, with the corresponding free circular limit.
Two independent GUE inputs can be combined as $(H_1^{(N)}+\mathrm{i}H_2^{(N)})/\sqrt{2}$ to produce a normalized complex Ginibre matrix, whose real and imaginary entry coordinates are independent centered Gaussians of variance $1/(2N)$.
Writing a general GUE polynomial in these variables introduces their adjoints, a case discussed in \cite[Question~1.8]{Han26}.
We prove almost-sure weak convergence for the full class of GUE polynomials.

\begin{theorem}\label{thm:brown}
Let $d\ge 1$ and fix $P\in\mathbb C\langle x_1,\ldots,x_d\rangle$, where $x_1,\ldots,x_d$ are self-adjoint indeterminates.
Let $H^{(N)}$ be as above, and let $s_1,\ldots,s_d\in\mathcal M$ be freely independent centered semicircular elements of variance one in a finite von Neumann algebra with faithful normal tracial state $\tau$.
Then the empirical measure
\[
 \mu_{P(H^{(N)})}:=\frac1N\sum_{k=1}^{N}\delta_{\lambda_k(P(H^{(N)}))}
\]
converges weakly almost surely to $\mu_{P(s)}^{\mathrm{Br}}$ for every coupling across sizes.
The eigenvalues are counted with algebraic multiplicity.
For each fixed $P$ and each coupling, the convergence holds on one event of probability one against every bounded continuous real-valued test function on $\mathbb C$.
\end{theorem}

We first state the estimate for a matrix pencil with identity constant term.
For a Hermitian tuple $H=(H_1,\ldots,H_d)$ of size $N$, let
\begin{equation}\label{eq:pencil}
 T_N(H):=I_m\otimes I_N+\sum_{j=1}^d R_j\otimes H_j,
\end{equation}
where the coefficient matrices $R_1,\ldots,R_d\in M_m(\mathbb C)$ are fixed.
We call such a pencil monic.
For a square matrix $B$, write $s_{\min}(B):=\min_{\|v\|=1}\|Bv\|$.

\begin{theorem}\label{thm:pencil}
Let $d,m\ge 1$, and let $R_1,\ldots,R_d\in M_m(\mathbb C)$ be arbitrary fixed matrices.
There exists $A>0$ such that
\begin{equation}\label{eq:pencil-floor}
 \sum_{N=1}^{\infty}\mathbb P\{s_{\min}(T_N(H^{(N)}))<N^{-A}\}<\infty .
\end{equation}
For each $M>0$ there exist $C,\alpha,\varepsilon_0>0$, independent of $N$, such that for $0<\varepsilon\le\varepsilon_0$,
\begin{equation}\label{eq:pencil-bound}
 \begin{split}
 &\mathbb P\{s_{\min}(T_N(H^{(N)}))\le\varepsilon,\
                \max_j\|H_j^{(N)}\|\le M\}\\
 &\hspace{2em}\le C N^{5/2}\varepsilon^\alpha
                  \exp(CN^{5/2}\varepsilon^\alpha).
 \end{split}
\end{equation}
\end{theorem}

For arbitrary fixed $A_0,\ldots,A_d\in M_m(\mathbb C)$, write
\begin{equation}\label{eq:affine-pencil}
 L_N(H):=A_0\otimes I_N+\sum_{j=1}^d A_j\otimes H_j,\qquad
 L:=A_0\otimes 1_{\mathcal M}+\sum_{j=1}^d A_j\otimes s_j .
\end{equation}
We use $L_N:=L_N(H^{(N)})$ for the random evaluation and $\tau_m:=\tr_m\otimes\tau$ for the normalized trace on $M_m(\mathcal M)$.

\begin{definition}\label{def:fk-determinant}
For $b\in M_m(\mathcal M)$ the Fuglede--Kadison determinant is $\Delta_{\tau_m}(b):=\exp(\tau_m(\log|b|))$, with $\log 0=-\infty$ and $\exp(-\infty)=0$ \cite[Definition~2.1(ii) and Remark~2.2(i)]{MS24}.
For a complex matrix $B$ of size $k$, write $\Delta_k(B):=|\det B|^{1/k}$ for the determinant with respect to the normalized trace \cite[Proposition~2.3(vi)]{MS24}.
\end{definition}

\begin{theorem}\label{thm:affine}
Let $L_N$ and $L$ be as in \eqref{eq:affine-pencil}, and put
\[
 U_N:=\frac{1}{mN}\log|\det L_N|,\qquad
 u:=\tau_m(\log|L|).
\]
For every coupling across sizes, $U_N\to u$ and $\Delta_{mN}(L_N)\to\Delta_{\tau_m}(L)$ almost surely.
For every $0<p<\infty$,
\[
 \mathbb E|\Delta_{mN}(L_N)-\Delta_{\tau_m}(L)|^p\longrightarrow 0.
\]
If $u>-\infty$, then $U_N$ is integrable for all sufficiently large $N$, and $\mathbb E|U_N-u|\to0$, in particular $\mathbb E U_N\to u$.
In this case there are $b,C,\varepsilon_0>0$ and $N_0$ such that
\begin{equation}\label{eq:global-floor}
 \mathbb P\{s_{\min}(L_N)\le\varepsilon\}\le C N^{5/2}\varepsilon^b
 \qquad
 \bigl(N\ge N_0,\ 0<\varepsilon\le\varepsilon_0,\ N^{5/2}\varepsilon^b\le1\bigr).
\end{equation}
If $u=-\infty$, then $U_N^+$ is integrable for every $N$ and $\mathbb E U_N\to-\infty$ in the extended real sense.
\end{theorem}

The first step is to show that a small singular value puts the Hermitian input tuple close to a tuple for which the pencil is singular.
In Section~\ref{sec:distance}, we compress an approximate kernel vector and the action of the inputs on its components to a space of dimension bounded in terms of $m$ and $d$.
A polynomial distance inequality in that dimension yields a nearby exact kernel equation, which we lift to the original space.
Section~\ref{sec:gaussian-tubes} then bounds the Gaussian measure of a neighborhood of the zero set of the determinant.
In Section~\ref{sec:affine}, a rank argument ensures that the affine determinant is a nonzero polynomial for all large sizes when $u$ is finite.
We then let the norm cutoff depend on the singular-value threshold to obtain \eqref{eq:global-floor}.
Gaussian concentration for truncated logarithms yields the almost-sure determinant limit.

In Section~\ref{sec:brown}, a Schur complement identifies the scalar logarithmic potential with the normalized logarithmic determinant of an affine pencil, up to the fixed pencil size.
Hermitization and logarithmic truncation identify the limit, following the method of \cite[Section~5]{Han26}.
Fubini's theorem and local square-integrability of the logarithmic kernel give pathwise local $L^1$ convergence of the potentials and hence weak convergence of the eigenvalue measures.
